%==========================================================================
%  TRES0312 Fysikk — Foiler-mal
%
%  MØrk/lys modus: kommenter ut én av de to linjene nedenfor.
%==========================================================================
\newif\ifdarkmode
\darkmodetrue      % ← mørk modus
%\darkmodefalse   % ← lys modus  (fjern % foran og legg % foran linjen over)

\documentclass[aspectratio=169, 12pt]{beamer}

\usepackage{amd-beamer}

%--------------------------------------------------------------------------
%  DOKUMENTINFO
%--------------------------------------------------------------------------
\title{Matematikk vi trenger}
\subtitle{TRES0312 --- Vektorer og modellering}
\author{Ben David Normann}
\date{\today}

%--------------------------------------------------------------------------
%  SEKSJONSKONTEKST
%  \setcounter{section}{N} setter gjeldende seksjonsnummer til N, slik at
%  \defn, \exm og \oppgave får samme nummerering som i kompendiet.
%  Tellerne nullstilles automatisk ved hvert \setcounter{section}{...}.
%  Bruk gjerne \setcounter{section}{N} rett før en frame der du bytter
%  seksjon, i stedet for én gang globalt.
%--------------------------------------------------------------------------
\setcounter{section}{2}   % Kapittel 2: "Matematiske verktøy vi trenger"

\begin{document}

%--------------------------------------------------------------------------
%  SLIDE 1 — TITTELSIDE
%--------------------------------------------------------------------------
\begin{frame}[plain]
  \titlepage
\end{frame}

%--------------------------------------------------------------------------
%  SLIDE 2 — Fem kjappe fra forrige gang
%--------------------------------------------------------------------------
\begin{frame}{}

  \repetisjon{Fem kjappe fra forrige gang}{%
    \begin{enumerate}[1.,itemsep=10pt]
      \item Hva kalles $a$ i en lineær funksjon $f(x) = b + ax$?
      \item Hva er $\cos(60^\circ)$?
      \item Skriv Pythagoras' setning for sidene i en rettvinklet trekant.
      \item Hva er en vektor?
      \item Hva er et koordinatsystem?
    \end{enumerate}
  }

\end{frame}

%--------------------------------------------------------------------------
%  SLIDE 3 — Definisjon: enhetsvektor
%--------------------------------------------------------------------------
\begin{frame}{Enhetsvektorer}

  \defn{Enhetsvektor}{
  En \textit{enhetsvektor} er en vektor med lengde 1 som angir en bestemt retning i rommet.
  }
  \pause
  \eksbox{
  I $xy$-planet finnes det to enhetsvektorer:
  \[
    \hat{\mathbf{x}} = [1,0], \qquad \hat{\mathbf{y}} = [0,1].
  \]
  Enhver vektor kan skrives som en kombinasjon av disse: $\vec{v} = [v_x,v_y] = v_x\hat{\mathbf{x}} + v_y\hat{\mathbf{y}}$.
  }

\end{frame}

%--------------------------------------------------------------------------
%  SLIDE 4 — Aktivitet: enhetsvektorer
%--------------------------------------------------------------------------
\begin{frame}{Aktivitet}

  \aktivitet{Vektorer som lineærkombinasjon av enhetsvektorer}{%
    \begin{enumerate}[a)]
      \item Skriv vektoren $\vec{A} = [4,\,-3]$ som en lineær kombinasjon av $\hat{\mathbf{x}}$ og $\hat{\mathbf{y}}$.
      \item Regn ut summen $\vec{u}+\vec{v}$ der $\vec{u} = 3\hat{\mathbf{x}} - 2\hat{\mathbf{y}}$ og $\vec{v} = -\hat{\mathbf{x}} + 5\hat{\mathbf{y}}$.
    \end{enumerate}
  }

\end{frame}

%--------------------------------------------------------------------------
%  SLIDE 5 — Dekomponering av vektorer
%--------------------------------------------------------------------------
\begin{frame}{Dekomponering av vektorer}

  Fysikkens lover formuleres med vektorer, uavhengig av koordinatsystem. Når vi skal regne, \textit{dekomponerer} vi vektorlikningen i hver retning:

  \graabox{%
    $\vec{v}(t) = \vec{v}_0 + \vec{a}t$
    \quad$\Longrightarrow$\quad
    $v_x(t) = v_{0x} + a_x t, \;\; v_y(t) = v_{0y} + a_y t$
  }

  \pause
  \merk{Fysikkens lover er uavhengig av koordinatsystem}{%
  Lengden til et hus er den samme enten du måler det med en tommestokk i meter eller i tommer.
  }

\end{frame}

%--------------------------------------------------------------------------
%  SLIDE 6 — Aktivitet: dekomponering av hastighet
%--------------------------------------------------------------------------
\begin{frame}{Aktivitet}

  \aktivitet{Dekomponering av hastighet}{%
  En ball kastes med fart $v = 20\,\mathrm{m/s}$ i en vinkel $\theta = 37^\circ$ over horisonten ($\cos 37^\circ \approx 0.80$, $\sin 37^\circ \approx 0.60$).
    \begin{enumerate}[a)]
      \item Finn komponentene $v_x$ og $v_y$.
      \item Verifiser at $|\vec{v}| = \sqrt{v_x^2+v_y^2} = 20\,\mathrm{m/s}$.
    \end{enumerate}
  }

\end{frame}

%--------------------------------------------------------------------------
%  SLIDE 7 — Aktivitet: trigonometri og vektorkomponenter
%--------------------------------------------------------------------------
\begin{frame}{Aktivitet}

  \aktivitet{Trigonometri og vektorkomponenter}{%
  En fugl flyr $10\,\mathrm{km}$ i retning N$60^\circ$Ø (det vil si $60^\circ$ øst for nord). Finn nord- og østkomponentene til forskyvningsvektoren. \textit{Hint: } vinkelen måles fra nordaksen, ikke østaksen.
  }

\end{frame}

%--------------------------------------------------------------------------
%  SLIDE 8 — Eksempel: posisjonsvektor i prosjektilbevegelse
%--------------------------------------------------------------------------
\begin{frame}{Eksempel: Kule saker}

  \begin{columns}[t]
    \begin{column}{0.52\textwidth}
      \vspace{0.4cm}
      \begin{center}
      \begin{tikzpicture}
      \node[anchor=south west, inner sep=0] (img) at (0,0)
        {\includegraphics[width=\linewidth]{Kapittel2:Kinematikk/Prosjektilbevegelse_malingsbombe_\ifdarkmode dark\else clean\fi.png}};
      \begin{scope}[x={(img.south east)-(img.south west)}, y={(img.north west)-(img.south west)}]
      \draw[-{Stealth[length=4mm,width=3mm]}, very thick, posGraphColor] (0.2646,0.1728) -- (0.3443,0.7188);
      \draw[-{Stealth[length=4mm,width=3mm]}, very thick, posGraphColor] (0.2646,0.1728) -- (0.5632,0.6313);
      \draw[-{Stealth[length=4mm,width=3mm]}, very thick, posGraphColor] (0.2646,0.1728) -- (0.6891,0.4577);
      \node[font=\footnotesize, posGraphColor] at (0.3727,0.5094) {$\vec{S}(t_1)$};
      \node[font=\footnotesize, posGraphColor] at (0.5044,0.4033) {$\vec{S}(t_2)$};
      \node[font=\footnotesize, posGraphColor] at (0.5588,0.2963) {$\vec{S}(t_3)$};
      \end{scope}
      \end{tikzpicture}
      \end{center}
    \end{column}
    \begin{column}{0.46\textwidth}
      \aktivitet{Skutt og sluppet kule}{%
      En paintball-kanonkule skytes horisontalt med farten $v_0 = 3\,\mathrm{m/s}$. Samtidig slippes en tilsvarende kule vertikalt. Se bort ifra luftmotstand.
      \begin{enumerate}[a)]
        \item Hvilken av dem når bakken først?
        \item Forsøk å bruke innsikten fra a) til å skrive ned uttrykk for $v_x(t)$ og $v_y(t)$.
      \end{enumerate}
      }
    \end{column}
  \end{columns}

\end{frame}

%--------------------------------------------------------------------------
%  SLIDE 9 — Eksempel: hastighetsvektor i prosjektilbevegelse
%--------------------------------------------------------------------------
\begin{frame}{Eksempel: Kule saker}

  \begin{columns}[t]
    \begin{column}{0.52\textwidth}
      \vspace{0.4cm}
      \begin{center}
      \begin{tikzpicture}
      \node[anchor=south west, inner sep=0] (img) at (0,0)
        {\includegraphics[width=\linewidth]{Kapittel2:Kinematikk/Prosjektilbevegelse_malingsbombe_\ifdarkmode dark\else clean\fi.png}};
      \begin{scope}[x={(img.south east)-(img.south west)}, y={(img.north west)-(img.south west)}]
      % Posisjonsvektorer
      \draw[-{Stealth[length=4mm,width=3mm]}, very thick, posGraphColor] (0.2646,0.1728) -- (0.3443,0.7188);
      \draw[-{Stealth[length=4mm,width=3mm]}, very thick, posGraphColor] (0.2646,0.1728) -- (0.5632,0.6313);
      \draw[-{Stealth[length=4mm,width=3mm]}, very thick, posGraphColor] (0.2646,0.1728) -- (0.6891,0.4577);
      \node[font=\footnotesize, posGraphColor] at (0.3727,0.5094) {$\vec{S}(t_1)$};
      \node[font=\footnotesize, posGraphColor] at (0.5044,0.4033) {$\vec{S}(t_2)$};
      \node[font=\footnotesize, posGraphColor] at (0.5588,0.2963) {$\vec{S}(t_3)$};
      % Hastighetsvektorer (samme horisontalkomponent $v_x$ hele veien, tangent til banen)
      \draw[-{Stealth[length=4mm,width=3mm]}, very thick, velGraphColor] (0.3443,0.7188) -- (0.4181,0.7109);
      \draw[-{Stealth[length=4mm,width=3mm]}, very thick, velGraphColor] (0.5632,0.6313) -- (0.6369,0.5663);
      \draw[-{Stealth[length=4mm,width=3mm]}, very thick, velGraphColor] (0.6891,0.4577) -- (0.7629,0.3072);
      \node[font=\footnotesize, velGraphColor] at (0.3835,0.7629) {$\vec{v}(t_1)$};
      \node[font=\footnotesize, velGraphColor] at (0.6162,0.6405) {$\vec{v}(t_2)$};
      \node[font=\footnotesize, velGraphColor] at (0.7517,0.4113) {$\vec{v}(t_3)$};
      % Hastighetsvektor for kula til venstre (samme høyde som $S(t_3)$, kun vertikal fart)
      \draw[-{Stealth[length=4mm,width=3mm]}, very thick, velGraphColor] (0.0627,0.4597) -- (0.0627,0.3092);
      \node[font=\footnotesize, velGraphColor] at (0.0898,0.3845) {$\vec{v}$};
      \end{scope}
      \end{tikzpicture}
      \end{center}
    \end{column}
    \begin{column}{0.46\textwidth}
      \aktivitet{Skutt og sluppet kule}{%
      En paintball-kanonkule skytes horisontalt med farten $v_0 = 3\,\mathrm{m/s}$. Samtidig slippes en tilsvarende kule vertikalt. Se bort ifra luftmotstand.
      \begin{enumerate}[a)]
        \item Hvilken av dem når bakken først?
        \item Forsøk å bruke innsikten fra a) til å skrive ned uttrykk for $v_x(t)$ og $v_y(t)$.
      \end{enumerate}
      }
    \end{column}
  \end{columns}

\end{frame}

%--------------------------------------------------------------------------
%  SLIDE 10 — Vektorlengde og -vinkel
%--------------------------------------------------------------------------
\begin{frame}{Størrelse og retning til en vektor}

  \defn{Vektorlengde}{
  Lengden til en vektor $\vec{v} = [v_x,v_y]$ er gitt ved
  \[
    v = |\vec{v}| = \sqrt{v_x^2 + v_y^2}.
  \]
  }
  \pause
  \defn{Vektorvinkel}{
  Vinkelen til en vektor $\vec{v} = [v_x,v_y]$, målt fra positiv $x$-retning, er gitt ved
  \[
    \theta = \arctan\!\left(\frac{v_y}{v_x}\right).
  \]
  }

\end{frame}

%--------------------------------------------------------------------------
%  SLIDE 11 — Aktivitet: størrelse og retning
%--------------------------------------------------------------------------
\begin{frame}{Aktivitet}

  \aktivitet{Størrelse og retning}{%
  Finn lengde og vinkel til vektoren $\vec{v} = [3,\,4]$.
  }

\end{frame}

%--------------------------------------------------------------------------
%  SLIDE 12 — Pause
%--------------------------------------------------------------------------
\begin{frame}{}
\begin{center}
\Huge{Pause!}
\end{center}
\end{frame}

%--------------------------------------------------------------------------
%  SLIDE 13 — Definisjon: vektoraddisjon
%--------------------------------------------------------------------------
\begin{frame}{Operasjoner med vektorer}

  \defn{Vektoraddisjon}{
  Addisjon av vektorer gjøres grafisk med hale-til-spiss-metoden. Algebraisk legges komponentene sammen retningsvis:
  \[
    \vec{v} + \vec{u} = [v_x+u_x, \; v_y+u_y].
  \]
  }

\end{frame}

%--------------------------------------------------------------------------
%  SLIDE 14 — Aktivitet: vektoraddisjon
%--------------------------------------------------------------------------
\begin{frame}{Aktivitet}

  \aktivitet{Vektoraddisjon}{%
  Gitt vektorene $\vec{B} = [3,\,1]$ og $\vec{C} = [1,\,2]$. Finn resultantvektoren $\vec{A} = \vec{B}+\vec{C}$ og beregn lengden $|\vec{A}|$.
  }

\end{frame}

%--------------------------------------------------------------------------
%  SLIDE 15 — Definisjon: skalarprodukt
%--------------------------------------------------------------------------
\begin{frame}{Skalarprodukt}

  \defn{Skalarprodukt}{
  Skalarproduktet av to vektorer $\vec{u}=[u_x,u_y]$ og $\vec{v}=[v_x,v_y]$ er
  \[
    \vec{u}\cdot\vec{v} = u_x v_x + u_y v_y = |\vec{u}|\,|\vec{v}|\,\cos\theta,
  \]
  der $\theta$ er vinkelen mellom vektorene.
  }

\end{frame}

%--------------------------------------------------------------------------
%  SLIDE 16 — Geometrisk tolkning
%--------------------------------------------------------------------------
\begin{frame}{}

  \begin{columns}[t]
    \begin{column}{0.52\textwidth}
      \merk{Geometrisk tolkning}{
      Skalarproduktet er et mål på hvor mye to vektorer peker i samme retning. Er $\theta=0^\circ$ er produktet maksimalt positivt, er $\theta=90^\circ$ er produktet null, og er $\theta=180^\circ$ er produktet maksimalt negativt.
      }
    \end{column}
    \begin{column}{0.46\textwidth}
      \vspace{0.3cm}
      \begin{center}
      \begin{tikzpicture}[scale=0.95]
        \draw[->, very thick, vecBcol] (0,0) -- (4.3,0) node[right, font=\small, vecBcol] {$\vec{B}$};
        \draw[->, very thick, vecAcol] (0,0) -- (2.4,1.7) node[above right, font=\small, vecAcol] {$\vec{A}$};
        \node[above left=1pt, font=\small, vecAcol] at (1.2,0.85) {$|\vec{A}|$};
        \draw[dashed, thick, textcolor!55] (2.4,1.7) -- (2.4,0);
        \draw[textcolor] (2.4,0) ++(-0.13,0) -- ++(0,0.13) -- ++(0.13,0);
        \draw[very thick, accGraphColor] (0,0) -- (2.4,0)
              node[midway, below=5pt, font=\small, textcolor] {$|\vec{A}|\cos\theta$};
        \draw[->, thick, textcolor] (0.65,0) arc (0:35:0.65);
        \node[font=\small, textcolor] at (0.92,0.20) {$\theta$};
      \end{tikzpicture}
      \end{center}
    \end{column}
  \end{columns}

\end{frame}

%--------------------------------------------------------------------------
%  SLIDE 17 — Aktivitet: skalarprodukt og vinkel
%--------------------------------------------------------------------------
\begin{frame}{Aktivitet}

  \aktivitet{Skalarprodukt og vinkel}{%
  Finn skalarproduktet av $\vec{u} = [3,\,4]$ og $\vec{v} = [1,\,-2]$, og bestem vinkelen mellom dem.
  }

\end{frame}

%--------------------------------------------------------------------------
%  SLIDE 18 — Definisjon: kryssprodukt
%--------------------------------------------------------------------------
\begin{frame}{Kryssprodukt}

  \begin{columns}[t]
    \begin{column}{0.52\textwidth}
      \defn{Kryssprodukt i 2D}{
      For $\vec{u}=[u_x,u_y]$ og $\vec{v}=[v_x,v_y]$ i $xy$-planet er kryssproduktets $z$-komponent
      \[
        (\vec{u}\times\vec{v})_z = u_x v_y - u_y v_x,
      \]
      med størrelse $|\vec{u}\times\vec{v}| = |\vec{u}|\,|\vec{v}|\,\sin\theta$.
      }
    \end{column}
    \begin{column}{0.46\textwidth}
      \vspace{0.3cm}
      \begin{center}
      \begin{tikzpicture}[scale=0.95]
        \fill[vecAcol!20!bodycolor] (0,0) -- (4.3,0) -- (6.7,1.7) -- (2.4,1.7) -- cycle;
        \draw[->, very thick, vecBcol] (0,0) -- (4.3,0) node[right, font=\small, vecBcol] {$\vec{B}$};
        \draw[->, very thick, vecAcol] (0,0) -- (2.4,1.7) node[above right, font=\small, vecAcol] {$\vec{A}$};
        \node[above left=1pt, font=\small, vecAcol] at (1.2,0.85) {$|\vec{A}|$};
        \draw[dashed, thick, textcolor!55] (2.4,1.7) -- (2.4,0);
        \draw[textcolor] (2.4,0) ++(-0.13,0) -- ++(0,0.13) -- ++(0.13,0);
        \draw[<->, thick, accGraphColor] (2.65,0) -- (2.65,1.7)
              node[midway, right=3pt, font=\small, textcolor] {$|\vec{A}|\sin\theta$};
        \draw[dashed, thick, textcolor!45] (4.3,0) -- (6.7,1.7);
        \draw[dashed, thick, textcolor!45] (2.4,1.7) -- (6.7,1.7);
        \node[font=\small, textcolor!75] at (4.6,1.3) {Areal $= |\vec{A}\times\vec{B}|$};
        \draw[->, thick, textcolor] (0.65,0) arc (0:35:0.65);
        \node[font=\small, textcolor] at (0.92,0.20) {$\theta$};
      \end{tikzpicture}
      \end{center}
    \end{column}
  \end{columns}

\end{frame}

%--------------------------------------------------------------------------
%  SLIDE 19 — Aktivitet: kryssprodukt
%--------------------------------------------------------------------------
\begin{frame}{Aktivitet}

  \aktivitet{Kryssprodukt}{%
  Finn kryssproduktet av $\vec{u} = [3,\,4]$ og $\vec{v} = [1,\,-2]$.
  }

\end{frame}

%--------------------------------------------------------------------------
%  SLIDE 20 — Matematisk modellering
%--------------------------------------------------------------------------
\begin{frame}{Matematisk modellering}

  \oppsum{Matematisk modellering}{%
    En matematisk modell av virkeligheten vil for oss inneholde:
    \begin{enumerate}[1.,itemsep=6pt]
      \item Representasjon av objektene involvert.
      \item Representasjon av påvirkningen mellom objektene (kreftene).
      \item Effektivt valg av koordinatsystem.
    \end{enumerate}
  }
  \pause
  \tanke{}{%
  Tegninger er ofte fysikernes beste venn --- det er ikke uten grunn at Feynman utarbeidet de såkalte Feynman-diagrammene.
  }

\end{frame}

%--------------------------------------------------------------------------
%  SLIDE 21 — Neste gang
%--------------------------------------------------------------------------
\begin{frame}
\begin{center}
  \Huge{Neste gang: Kapittel 3 --- Kinematikk!}
  \begin{itemize}
    \item Hva er kinematikk?
    \item Strekning og hastighet
    \item Akselerasjon
    \item Bevegelseslikningene
  \end{itemize}
\end{center}
\end{frame}


\end{document}
